\section{Training on public CT projections with a twin regulariser}
\label{sec:method_v2}
% PLAN (about 1.5 pages). This section is new text for the v2 recipe; methods.tex (v1 recipe, physics, audit definition) is untouched.
%   4.1 Qure projections as an exactly additive source: what a row contains, why the bone image is exact (composite = a*soft + b*bone + c*lung + d),
%       the augmentation order of s4/qp.py, and why monotone tone curves keep the decomposition exact.
%   4.2 Twin regulariser: definition, loss, the weight w, and the variants that form the ablation (stop-gradient, support size, whole image, number
%       of pasted lesions, insertion without the tie term). Relation to the v1 counterfactual insertion (CLIA) in methods.tex.
%   4.3 Ensemble and uncertainty: K=3 members, mean bone image, spread, split-conformal scale, what is calibrated where (SAP_v2.md, H5').
%   4.4 Training budget and what is identical across arms (same optimiser and schedule as methods.tex).
% Verify every formula below against s4/train.py, s4/qp.py and s4/uncertainty.py at the freeze commit before submission.
\subsection{Qure projections as an exactly additive source}
Each row of the Qure.ai data holds a composite radiograph $T$ and the translated bone, soft-tissue and lung projections $B_t$, $S_t$ and $L$
(8-bit, $512\times512$, one sample per CT volume at 100~kVp). We verified that $T=a\,S_t+b\,B_t+c\,L+d$ holds up to rounding, with a per-sample bone
weight $b$ (\path{tools/qp_prep.py}); the additive bone image of the normalised input $X=T/255$ is therefore $B=b\,B_t/255$ and the remainder is
$S=X-B$. Over all \num{21887} rows the least-squares fit leaves a mean absolute residual of 0.36 grey levels per projection (largest 3.2 of 255).
The rows are split by patient into \num{21051} training, 406 validation and 430 test projections.

Training samples are composed as in Section~\ref{sec:methods}: bone re-mixing ($X'=S+wB$, exact), one random affine transform shared by all maps,
lesion insertion into $S$ with the bone target untouched, a monotone gamma tone curve ($X'=g(S+B)$, $S'=g(S)$, $B'=X'-S'$), linear blur applied to
$X'$ and $B'$, and noise and scatter applied to $X'$ only. Defining $B'=X'-S'$ keeps the decomposition exact under any monotone tone curve.

\subsection{Twin lesion-invariance regulariser}
Let $f$ be the suppressor, $x$ a training image and $b$ its exact bone image. A \emph{twin} of $x$ is $x^{+}=\min(x+n,1)$, where $n\ge0$ is the
signal of a few lesion patches taken from the lesion library of \emph{training} CTs and pasted at random positions in the lung fields, in composite
units and after all augmentation. The bone target of $x^{+}$ is the same $b$. Let $M$ be the lesion support dilated by $3$~px. The loss is
\begin{equation}
\begin{gathered}
\mathcal{L}=\tfrac12\big[\mathcal{L}_{\mathrm{sup}}(f(x),b)+\mathcal{L}_{\mathrm{sup}}(f(x^{+}),b)\big]+w\,\mathcal{L}_{\mathrm{inv}},\\
\mathcal{L}_{\mathrm{inv}}=\frac{\sum_p M_p\,\big|f(x^{+})_p-f(x)_p\big|}{\max\big(1,\sum_p M_p\big)},
\end{gathered}
\label{eq:twin}
\end{equation}
where $\mathcal{L}_{\mathrm{sup}}$ is the suppression loss of Section~\ref{sec:methods} and $w$ is the twin weight. Half of every batch consists of twins, so
the cost per step does not change. The supervised term already asks for the unchanged bone image on $x^{+}$; $\mathcal{L}_{\mathrm{inv}}$ ties the
two predictions directly on the lesion support. Each twin carries between one and two pasted lesions (uniform), each with a random scale and
amplitude; the supervised term is averaged over all base images and twins of the batch, and $\mathcal{L}_{\mathrm{inv}}$ is normalised per image by
its support and then averaged over the twins.

The ablation (Section~\ref{sec:experiments}) changes one part at a time: the twin weight $w$; a stop-gradient on $f(x)$; a support dilated further with a
$31\times31$ window, or the whole image, instead of $M$; up to four instead of up to two pasted lesions; and lesion insertion into the base image
without the tie term (the counterfactual insertion of Section~\ref{sec:methods}). All of these arms are in the ablation of
Section~\ref{sec:res_twin} (development data); the twin weights form the confirmatory dose-response.

\subsection{Deep ensemble and uncertainty}
The reported model is a deep ensemble of $K=3$ suppressors trained with seeds $0$, $1$ and $2$ \citep{lakshminarayanan2017ensembles}; seed $3$ is a reserve.
The bone image is the ensemble mean and the soft image is $x$ minus that mean. The ensemble standard deviation is the per-pixel spread. A split-conformal
scale $q$ is the $\lceil (n+1)(1-\alpha)\rceil/n$ quantile of $|b-\bar f|/(\sigma+\varepsilon)$ pooled over calibration pixels of the Qure-projection
validation CTs, so that $[\bar f-q(\sigma+\varepsilon),\,\bar f+q(\sigma+\varepsilon)]$ has nominal coverage $1-\alpha=0.9$
\citep{lei2018conformal,angelopoulos2023conformal}. Calibration treats pixels as exchangeable, so the guarantee is marginal over pixels, not
over images. The sparsification curve is the mean absolute error of the pixels that remain after removing the fraction $f\in[0,0.95]$ with the
largest spread; AUSE is the area between this curve and the oracle curve that removes pixels by their error (0 is perfect).

\subsection{Training budget}
The suppressor keeps the nested decoder and CBAM attention of Section~\ref{sec:methods} \citep{zhou2018unetpp,woo2018cbam}. The dose-response and
ablation arms use the SwinV2-T encoder of that section \citep{liu2022swinv2}; ours uses a ConvNeXt-V2-T encoder \citep{woo2023convnextv2}
(FCMAE pre-training, ImageNet-22k then ImageNet-1k fine-tuning \citep{deng2009imagenet}; \texttt{timm} \citep{wightman2019timm}
\path{convnextv2_tiny.fcmae_ft_in22k_in1k}) with a learning
rate $0.25\times$ that of the decoder. Every v2 run is trained at $512$~px with AdamW \citep{loshchilov2019adamw} (learning rate $2\times10^{-4}$,
batch 8), the schedule and loss of Section~\ref{sec:methods}, and \num{300000} samples (\num{37500} steps); evaluation uses the EMA weights. The
encoder screen uses \num{100000} samples (Section~\ref{sec:experiments}). The seeds per arm are listed in the generated tables.
